\section{An endomorphism-stable weakly admissible filtration}
\label{sec:filtration}

We construct a filtration on first crystalline cohomology that preserves an
additional divisor.  The construction takes place entirely on the given
isocrystal.  Its realization by an abelian lift will be a separate step.
Generic transversality is a basic ingredient of the admissible-filtration
existence theorems of Fontaine--Rapoport
\cite[Theorem~1 and Lemmas~1.1--1.2]{FontaineRapoport2005}.
We impose the additional endomorphism and polarization constraints
explicitly. The bounded-degree argument below chooses the required
coordinates over a finite algebraic extension, while testing every
subisocrystal over the fixed coefficient field.

Throughout this section, $\Kzero=\operatorname{Frac}W(\overline{\F}_p)$,
and valuations are normalized by $v(p)=1$.  If $D$ is an isocrystal over
$\Kzero$ and $K/\Kzero$ is finite, a subspace $L\subset D_K$ defines the
one-step filtration
\[
 \Fil^iD_K=
 \begin{cases}
 D_K,&i\leq 0,\\
 L,&i=1,\\
 0,&i\geq 2.
 \end{cases}
\]
For an isocrystal subobject $P\subset D$, its Hodge number for this
filtration is $t_H(P)=\dim_K(P_K\cap L)$.  Its Newton number $t_N(P)$ is
the sum of its slopes with multiplicities.  Weak admissibility means that
$t_H(D)=t_N(D)$ and $t_H(P)\leq t_N(P)$ for every isocrystal subobject
$P$ over $\Kzero$.

Finite local-field descent is understood as in Section~\ref{sec:transport}:
the isocrystal and its indicated tensors descend to the maximal unramified
subfield of a finite extension of \(\Q_p\), and the filtration descends
to that finite extension.

\begin{theorem}[The filtration construction]
\label{thm:filtration}
Let $A_0/\overline{\F}_p$ be an abelian variety equipped with a rational
action of a CM field $E$ and a polarization whose Rosati involution
restricts to conjugation on $E$.  Suppose that
\[
 D=H^1_{\mathrm{cris}}(A_0/W(\overline{\F}_p))[1/p]
\]
has rank $2s$ over $E\otimes_\Q\Kzero$, where $s\geq 1$.  Suppose also
that there is an invertible $u\in\End^0(A_0)$ such that
\[
 u^\dagger=u,\qquad ue=\bar e u\quad(e\in E).
\]
Let $\beta$ be the dual polarization pairing on $D$, with values in
the isocrystal line of slope $1$.

Then there are a finite extension $K/\Kzero$ and a subspace $L\subset
D_K$ with the following properties:
\begin{enumerate}
\item The one-step filtration with $\Fil^1D_K=L$ is weakly admissible.
\item $L$ is stable under the cohomological actions of $E$ and $u$,
and is Lagrangian for $\beta$.
\item After extending to an algebraically closed field $\Omega/K$,
one has $\dim_\Omega L_\tau=s$ for every embedding
$\tau:E\hookrightarrow\Omega$.
\item The isocrystal, its indicated tensors, and this filtration have
finite local field descent.
\end{enumerate}
Moreover, geometrically $L$ belongs to the balanced polarized
cocharacter orbit.  More explicitly, define
\[
 G_D=\left\{g\in\GSp(D,\beta)\cap
 \GL_{E\otimes_\Q\Kzero}(D):
 \det_{E\otimes_\Q\Kzero}(g)=c(g)^s\right\},
\]
where the scalar on the right is viewed in
$(E\otimes_\Q\Kzero)^\times$.  The balanced $E$-stable
$\beta$-Lagrangians form a single $G_D(\Omega)$-orbit.  Each is the
weight-$1$ space of a cocharacter of $G_{D,\Omega}$ with weights $0,1$.
In particular $L$ belongs to the orbit of any original filtration of
this type under a tensor-preserving comparison.
\end{theorem}

The proof will construct \(L\) from Lagrangians in finitely many
symplectic spaces. The two lemmas below supply the required simultaneous
choice: minimum intersection with a subspace is ensured by a nonzero
minor of bounded degree, and all such minors over a fixed coefficient
field can be made nonzero over one finite extension.

\begin{lemma}[Bounded-degree avoidance over a finite extension]
\label{lem:avoidance}
Let $K_1/\Kzero$ be a finite extension, and let $r,b\geq 1$ be integers.
There exist an element $\theta$ algebraic over $\Q_p$ and a point
\[
 (x_0,\ldots,x_{r-1})\in K_1(\theta)^r
\]
at which every nonzero polynomial in $K_1[X_0,\ldots,X_{r-1}]$ of total
degree at most $b$ is nonzero.  The coordinates $x_i$ can all be chosen
in $\Q_p(\theta)$.
\end{lemma}

\begin{proof}
Choose an integer $N>b$ and put
\[
 B=b\sum_{i=0}^{r-1}N^i.
\]
Let $e$ be the ramification index of $K_1/\Kzero$, so that
$v(K_1^\times)=\frac1e\Z$.  Choose a prime $q>\max(B,e)$, and choose
$\theta$ with $\theta^q=p$ in an algebraic closure of $K_1$.  The value
$v(\theta)=1/q$ has order $q$ modulo $\frac1e\Z$.  Consequently the
ramification index, and hence the degree, of $K_1(\theta)/K_1$ is at
least $q$.  The polynomial $X^q-p$ gives the reverse inequality, so
\[
 [K_1(\theta):K_1]=q>B.
\]

Set $x_i=\theta^{N^i}$.  Under the substitution $X_i=T^{N^i}$,
distinct monomials of total degree at most $b$ have distinct exponents:
each exponent of an individual variable is smaller than $N$, and
uniqueness follows from base-$N$ expansion.  Thus a nonzero polynomial
of total degree at most $b$ becomes a nonzero polynomial in $T$ of
degree at most $B$.  It cannot vanish at $\theta$, whose degree over
$K_1$ exceeds $B$.
\end{proof}

\begin{remark}
The coefficient field in Lemma~\ref{lem:avoidance} is the fixed field
$K_1$.  The conclusion concerns all bounded-degree polynomials over that
field, with no restriction on their cardinality.  It would be false with
$K_1(\theta)$ in place of $K_1$ as the coefficient field.
\end{remark}

\begin{lemma}[Intersections with Lagrangians]
\label{lem:lagrangian-intersection}
Let $M$ be a symplectic space of dimension $2m$ over a field $k$ of
characteristic zero, and let $W\subset M$ have dimension $a$.  A nonempty
open subset of the Lagrangian Grassmannian consists of subspaces $L$
such that
\[
 \dim(L\cap W)=\max(0,a-m).
\]
On any fixed symplectic affine big cell, at least one nonzero polynomial
of degree at most $m$ has nonvanishing locus contained in this open
subset.  This degree bound is independent of $W$.
\end{lemma}

\begin{proof}
We first work over an algebraic closure.  The Lagrangian Grassmannian
is geometrically integral of dimension $m(m+1)/2$; its affine big cells
are parametrized by symmetric $m\times m$ matrices.  If $1\leq a\leq m$,
consider the projective incidence variety of pairs $(\ell,L)$ with
$\ell\subset W$ a line and $\ell\subset L$ a Lagrangian.  The fiber over
$\ell$ is the Lagrangian Grassmannian of the symplectic space
$\ell^\perp/\ell$, and therefore the incidence variety has dimension
\[
 a-1+\frac{m(m-1)}2<\frac{m(m+1)}2.
\]
Its image is a proper closed subset.  Its complement consists precisely
of the Lagrangians missing $W$.  The case $a=0$ is immediate.  For
$a>m$, use
\[
 \dim(L\cap W)=a-m+\dim(L\cap W^\perp)
\]
and apply the preceding argument to $W^\perp$.

For the assertion on equations, choose coordinates on the fixed big
cell in which $L$ is the column space of
\[
 \begin{pmatrix}I_m\\ T\end{pmatrix},\qquad T=T^{\mathsf t}.
\]
The minimum-intersection condition says that the composite of this
matrix with the quotient map $M\longrightarrow M/W$ has rank
$\min(m,2m-a)$.  Its entries are affine linear polynomials in the
independent entries of $T$.  The nonempty open subset already proved
meets this big cell, so some minor of the required size is a nonzero
polynomial over $k$.  Its degree is at most $m$, and its nonvanishing
implies the required rank.  If the required rank is zero, use the
constant polynomial $1$.
\end{proof}


The two lemmas reduce simultaneous transversality to a uniform bound on
the degrees of the relevant minors. We now identify the symplectic
spaces on which to apply them and the inequalities they must enforce.

\begin{proof}[Proof of Theorem~\ref{thm:filtration}]
We write $e$ and $u$ also for their cohomological actions on $D$.
These actions commute with Frobenius and satisfy
\begin{equation}
\label{eq:filtration-adjoints}
 \beta(ux,y)=\beta(x,uy),\qquad
 \beta(ex,y)=\beta(x,\bar e y),\qquad ue=\bar e u.
\end{equation}
There is no change to the last relation from contravariance.  Indeed,
pulling back $ue=\bar e u$ gives $e^*u^*=u^*\bar e^*$; replacing $e$
by $\bar e$ gives $u^*e^*=\bar e^*u^*$.  Likewise, if $J$ is the
homological polarization matrix, $U^{\mathsf t}J=JU$ implies
$UJ^{-1}=J^{-1}U^{\mathsf t}$, which is the self-adjointness of the
cohomological action $U^{\mathsf t}$ for the dual pairing.  The
Rosati-conjugate adjoint identity for $E$ follows in the same way.

The pullback $U^{\mathsf t}$ of an individual endomorphism is
distinct from the contragredient group action $g\mapsto g^{-\mathsf t}$.
If a homological similitude $g$ has multiplier $c_V$ and satisfies
$\det_E(g)=c_V^s$, its dual action has multiplier $c_D=c_V^{-1}$
and satisfies $\det_E(g^{-\mathsf t})=c_D^s$.
Thus $G_D$ is the dual form of the homological group defined by this
determinant condition, with its cohomological multiplier.

\medskip\noindent
\emph{Reduction to the stable subobjects.}
First consider any $E,u$-stable candidate $L\subset D_K$, where
$K/\Kzero$ is finite. We show that the weak-admissibility inequalities
need only be tested on $E,u$-stable subobjects. This resembles the use of canonical
destabilizing objects in \cite[Section~2]{FontaineRapoport2005},
but we use total defect, not defect divided by dimension, and give
the required elementary argument.
For any isocrystal subobject $P\subset D$ over $\Kzero$, set
\[
 f(P)=t_H(P)-t_N(P).
\]
These values form a finite set: the Hodge number is a bounded integer,
and the Newton number is determined by bounded integer multiplicities
among the finitely many slopes of $D$.  Let $a$ be their maximum.
Newton numbers are additive in short exact sequences, while
\[
 (P_K\cap L)+(Q_K\cap L)\subset (P+Q)_K\cap L.
\]
Consequently
\begin{equation}
\label{eq:filtration-supermodularity}
 f(P)+f(Q)\leq f(P\cap Q)+f(P+Q).
\end{equation}
If $P$ and $Q$ both attain $a$, the right-hand terms are each at
most $a$, so both attain $a$.  Thus maximizers are closed under
intersection and sum.  Choose a maximizer $P_{\max}$ of largest
dimension.  Adding any other maximizer to it shows that the other
one is contained in it.  Hence $P_{\max}$ is the unique largest
maximizer.

The operators $u$ and $e\in E^\times$ are filtration-preserving
isocrystal automorphisms.  They preserve $f$ and hence fix
$P_{\max}$, by its uniqueness.  Thus $P_{\max}$ is $E,u$-stable.
Consequently, if $t_H(D')\leq t_N(D')$ for every $E,u$-stable
subobject $D'$, then $a=f(P_{\max})\leq0$. All the
weak-admissibility inequalities follow. We must therefore construct
an $E,u$-stable Lagrangian with these restricted inequalities and
$t_H(D)=t_N(D)$.

\medskip\noindent
\emph{Semisimplicity and symplectic blocks.}
The Rosati trace form
\[
 \langle a,b\rangle=\Tr(ab^\dagger)
 \quad\text{on }\End^0(A_0)\otimes_\Q\R
\]
is positive definite \cite[Theorem~17.3]{MilneAV}.  Left multiplication
by $u$ is self-adjoint for this form, since trace cyclicity gives
\[
 \langle ua,b\rangle=\Tr(uab^\dagger)
 =\Tr(ab^\dagger u)=\langle a,ub\rangle.
\]
It is therefore diagonalizable.  Its minimal polynomial equals the
minimal polynomial of $u$ in the endomorphism algebra, as is seen by
applying a polynomial in the multiplication operator to the identity.
Thus $u$, and hence $u^2$, acts semisimply on $D$.  The relation in
\eqref{eq:filtration-adjoints} also shows that $u^2$ commutes with $E$.

Put
\[
 B(x,y)=\beta(x,uy).
\]
This is a nondegenerate alternating pairing.  Indeed, invertibility is
immediate, and
$B(y,x)=-\beta(ux,y)=-B(x,y)$.  In addition,
\[
 B(ex,y)=\beta(x,\bar euy)=\beta(x,uey)=B(x,ey),
 \qquad B(u^2x,y)=B(x,u^2y).
\]
After a finite splitting extension $K_1/\Kzero$, write
\[
 D_{K_1}=\bigoplus_{\tau,z}M_{\tau,z},
 \qquad
 M_{\tau,z}=\{x:ex=\tau(e)x\ (e\in E),\ u^2x=zx\}.
\]
We omit zero summands.  The self-adjoint identities for $B$ show that
distinct summands are orthogonal.  Its restriction to each summand is
therefore nondegenerate.  Write $\dim M_{\tau,z}=2m_{\tau,z}$.

For $\bar\tau=\tau\circ\overline{(\cdot)}$, the operator $u$ maps
$M_{\tau,z}$ isomorphically to $M_{\bar\tau,z}$.  On these spaces,
\begin{equation}
\label{eq:filtration-u-similitude}
 B(ux,uy)=zB(x,y).
\end{equation}
Here $z\neq 0$, and $\tau\neq\bar\tau$.  Choose a $B$-Lagrangian
$L_{\tau,z}$ in one member of each opposite pair, and prescribe
\begin{equation}
\label{eq:filtration-opposite-choice}
 L_{\bar\tau,z}=uL_{\tau,z}.
\end{equation}
Equation~\eqref{eq:filtration-u-similitude} shows that the latter is
Lagrangian as well.  Their direct sum $L$ is $E$-stable and $u$-stable;
on applying $u$ twice, one obtains $zL_{\tau,z}=L_{\tau,z}$.  Since
\[
 \beta(x,y)=B(x,u^{-1}y),
\]
$L$ is also $\beta$-Lagrangian.  Its dimension at each $E$-embedding
is $s$.  It remains to make the independent choices appropriately.

\medskip\noindent
\emph{The inherited Newton grading.}
Use the canonical slope decomposition
\[
 D=\bigoplus_\lambda D_\lambda.
\]
The isocrystal category over $\overline{\F}_p$ is semisimple, and
this decomposition is functorial
\cite[\S3.2, Proposition~3.3]{ColmezFontaine2000}.  The slopes of $D$
belong to $[0,1]$: its integral Dieudonn\'e operators \(F,V\)
satisfy \(FV=p\), so both the slopes of \(F\) and those of \(V\)
are nonnegative.  All the endomorphisms above preserve every
$D_\lambda$, and the projectors onto the simultaneous eigenspaces
preserve its scalar extension.  Thus each block has a grading
\[
 M_{\tau,z}=\bigoplus_\lambda M_{\tau,z,\lambda},\qquad
 M_{\tau,z,\lambda}=M_{\tau,z}\cap(D_\lambda)_{K_1}.
\]
We use this grading even if Frobenius permutes the labels $(\tau,z)$.
No extension of Frobenius to $K_1$, and no isocrystal structure on an
individual block, is required.

Before scalar extension, $B$ is a pairing of isocrystals into the
slope-$1$ line.  Hence $B(D_\lambda,D_\mu)=0$ unless
$\lambda+\mu=1$.  Nondegeneracy on an individual block $M$ then gives
perfect pairings between $M_\lambda$ and $M_{1-\lambda}$.  In
particular, if $\dim M=2m$,
\begin{equation}
\label{eq:filtration-block-newton}
 \sum_\lambda\lambda\dim M_\lambda=m.
\end{equation}

Let $D'\subset D$ be an isocrystal subobject over $\Kzero$ stable
under $E$ and $u$.  Its intersections $M'=D'_{K_1}\cap M$ decompose
$D'_{K_1}$ and inherit the same slope grading.  Put $a=\dim M'$.
The slope sum of the graded quotient $M/M'$ is at most $2m-a$,
because all its slopes belong to $[0,1]$.  By
\eqref{eq:filtration-block-newton},
\begin{equation}
\label{eq:filtration-block-bound}
 \sum_\lambda\lambda\dim M'_\lambda
 \geq \max(0,a-m).
\end{equation}
The additional bound by zero follows from the nonnegative slopes.
The right side is also the smallest possible dimension of the
intersection of a Lagrangian in $M$ with $M'$.
Thus attaining this minimum in every block will give the
weak-admissibility inequalities for all the stable subobjects.

\medskip\noindent
\emph{Simultaneous choices and finite local field descent.}
We may arrange the splitting field and coordinates used above as
follows.  Choose a finite field $k_0\subset\overline{\F}_p$ over
which $A_0$, its polarization, $u$, and a rational basis of $E$ are
defined.  The underlying isocrystal and the indicated tensors then
descend to $\operatorname{Frac}W(k_0)$.  Choose a finite local field
$H/\Q_p$ containing this field and splitting the actions of $E$ and
$u^2$.  The eigenspaces and their symplectic bases can be constructed
over $H$ by linear algebra.  We take $K_1=\Kzero H$.  In particular,
the independent Lagrangian Grassmannians, fixed symplectic big cells
in them, and the maps induced by $u$ all have coordinates over $H$.

Take the product of these big cells, one for each independent choice
in \eqref{eq:filtration-opposite-choice}, and let $r$ be its number
of affine coordinates.  For every $E,u$-stable isocrystal subobject
$D'$ over $\Kzero$ and each independent block $M$, the space
$M'=D'_{K_1}\cap M$ is defined over the \emph{fixed} field $K_1$.
Lemma~\ref{lem:lagrangian-intersection} supplies a nonzero polynomial
over $K_1$, of degree at most $m$, whose nonvanishing ensures
\begin{equation}
\label{eq:filtration-minimal-intersection}
 \dim(L_{\tau,z}\cap M')=\max(0,\dim M'-m).
\end{equation}
We regard this as a polynomial on the whole product by ignoring the
other coordinates.  A single bound $b=\max m_{\tau,z}$ works for all
these polynomials, independently of $D'$.  On the opposite block,
the same intersection equality follows from $u$-stability of $D'$
and \eqref{eq:filtration-opposite-choice}.

Apply Lemma~\ref{lem:avoidance} to $K_1,r,b$.  Its one point satisfies
\eqref{eq:filtration-minimal-intersection} for every one of the
subobjects just specified, simultaneously.  This use of the lemma
does not assume that the collection of subobjects is countable or
that its members descend to $H$.  Their definition over $K_1$ and
the uniform polynomial degree bound are the only requirements.
There are nonzero blocks with $m\geq1$, so $r,b\geq1$ as required.

Put $K=K_1(\theta)=\Kzero H(\theta)$.  The resulting Lagrangian is
defined over $H(\theta)$ in the chosen coordinates, and $H(\theta)$
is finite over $\Q_p$.  Extending the original isocrystal to the
maximal unramified subfield of $H(\theta)$ therefore gives finite
local field descent of the complete filtered data.  In particular,
the construction has not obtained its generic point by introducing
transcendental parameters.

Summing \eqref{eq:filtration-minimal-intersection} and
\eqref{eq:filtration-block-bound} over all blocks gives
\begin{equation}
\label{eq:filtration-stable-test}
 t_H(D')\leq t_N(D')
 \quad\text{for every $E,u$-stable isocrystal subobject $D'\subset D$.}
\end{equation}
For $D$ itself, both numbers equal $\frac12\dim D$ by
\eqref{eq:filtration-block-newton}. The reduction to stable subobjects
proved at the start now gives all the weak-admissibility inequalities
for subobjects of $D$ over $\Kzero$. Each subobject over the smaller unramified
field of descent extends to one of these subobjects. Hence the
descended filtered isocrystal is weakly admissible as well.

\medskip\noindent
\emph{The geometric cocharacter orbit.}
Over $\Omega$, the polarization pairs $D_\tau$ only with
$D_{\bar\tau}$.  A balanced $E$-stable Lagrangian is determined by an
$s$-dimensional subspace of $D_\tau$ for one member of each pair;
the subspace in the opposite member is its annihilator.  The
subgroup of $G_D$ with multiplier $1$ contains the product of
$\SL(D_\tau)$ acting on opposite spaces by the inverse dual
representations.  These special linear groups are transitive on
the $s$-dimensional subspaces.  Thus all the indicated Lagrangians
form one $G_D(\Omega)$-orbit.

Finally choose an $E$-stable $\beta$-Lagrangian complement to one
such $L$.  Such a complement is obtained separately on each pair
of opposite embedding spaces.  Acting by $t$ on $L$ and by $1$
on the complement gives a cocharacter with multiplier $t$ and
$E$-determinant $t^s$ at each embedding.  It lies in $G_D$ and has
the stated filtration.  This proves the final assertions.
\end{proof}

\begin{corollary}
\label{cor:filtration-admissible}
The filtration of Theorem~\ref{thm:filtration}, with its finite local
field descent, is admissible: the descended filtered isocrystal is
the filtered crystalline module of a crystalline representation.
\end{corollary}

\begin{proof}
Apply the weakly-admissible-implies-admissible theorem of
Colmez--Fontaine with monodromy operator zero
\cite[Theorem~A]{ColmezFontaine2000}.
\end{proof}

\begin{corollary}[The divisor filtration condition]
\label{cor:filtration-divisor}
Let \(D,L\) be as in Theorem~\ref{thm:filtration}, and let \(\psi\)
and \(U\) denote the homological polarization and the homological
action of \(u\) on \(D_K^\vee\). The untwisted divisor form
\(\psi(\cdot,U\cdot)\) belongs to
\[
 \Fil^1\bigl(\bigwedge^2D_K\bigr)=L\wedge D_K,
\]
and hence to filtration degree zero after its codimension-one Tate
twist. If a marked power is realized by an abelian lift, the
corresponding pulled-back divisor satisfies the filtration condition
for its symmetric homomorphism in Lemma~\ref{lem:lift-hom}.
\end{corollary}

\begin{proof}
Let \(H\subset D_K^\vee\) be the annihilator of \(L\). It is
Lagrangian for \(\psi\) and stable under \(U\): for \(a\in L\)
and \(h\in H\),
\[
 a(Uh)=(u^*a)(h)=0.
\]
Therefore \(\psi(h,Uh')=0\) for \(h,h'\in H\). The divisor form
thus lies in
\[
 \ker\bigl(\bigwedge^2D_K\longrightarrow\bigwedge^2H^\vee\bigr)
       =L\wedge D_K.
\]
This proves the assertion on the fixed filtered vector space. After
realization of a marked power, pullback respects this filtration
condition, which is precisely the condition on the associated
symmetric homomorphism in Lemma~\ref{lem:lift-hom}.
\end{proof}
